\documentclass[11pt,a4paper]{article}
\usepackage[margin=2.2cm]{geometry}
\usepackage[T1]{fontenc}
\usepackage{lmodern,microtype,booktabs,longtable,array,amsmath,xcolor}
\usepackage[colorlinks=true,linkcolor=blue!50!black,citecolor=blue!50!black,urlcolor=blue!50!black]{hyperref}
\newcommand{\sw}[1]{\texttt{#1}}
\newcommand{\code}[1]{\texttt{#1}}
\setlength{\parskip}{4pt}\setlength{\parindent}{0pt}
\title{Open-charm extension of AMPT v1.26t9b/v2.26t9b:\\ what was changed, where, and why}
\author{Reetanshu K. Pandey\\ \small Dr B R Ambedkar National Institute of Technology Jalandhar}
\date{Release accompanying the repository \code{ampt-charm}, \today}
\begin{document}
\maketitle

\begin{abstract}
This note lists every change made to the public AMPT model v1.26t9b/v2.26t9b~\cite{ampt2005,lin2021} to study open charm
in Au+Au collisions at $\sqrt{s_{NN}}=200$ GeV. For each change it gives the physics motivation and the paper it
follows, the file, routine and source lines that implement it, the switch that turns it on, and the check that it
works. All changes are either a patch to the seven stock source files (170 lines added, 17 changed) or four new files
(\code{charm\_mod.f}, \code{coal\_new.f}, \code{hqctl.inc}, \code{eps09s\_f.cpp}). Every change is behind an
environment switch that defaults to off, so with no switch set the program behaves as stock AMPT. The appendix tables
are generated from the repository (\code{docs/tools/code\_map.py}), so the line numbers are those of the patched tree
that \code{scripts/build.sh} creates in \code{build/ampt-src/}.
\end{abstract}

\tableofcontents

\section{How to inspect the changes yourself}
\begin{enumerate}
\item \code{scripts/setup\_deps.sh} downloads the official AMPT archive from Z.-W.~Lin's web page and checks its
  SHA-256 (\code{9441b5f7\ldots}). The seven Fortran files of that archive are the base; no other version is used.
\item \code{scripts/build.sh} unpacks it, applies \code{patches/ampt-v1.26t9b-v2.26t9b\_charm.patch} and copies the
  four new files of \code{src\_new/}. Read the patch directly: it has 46 hunks (Table~\ref{tab:hunks}), each opening
  with a comment that starts with \code{cdmeson}. To list them: \code{grep -n cdmeson build/ampt-src/*.f}.
\item With no \code{AMPT\_*} variable set, every new routine returns at once and every stock branch is taken. This was
  checked byte for byte on \code{ana/ampt.dat} against the unmodified program (gate G0 of the development log).
  \code{configs/stock.env} runs the program this way.
\item \code{tests/test\_exact\_reproduction.sh} regenerates events of the paper production from their seeds and
  compares them with the files written on the batch farm (Sec.~\ref{sec:verify}).
\end{enumerate}

\section{Event flow and where the charm code enters}
The AMPT string-melting chain~\cite{ampt2005} is unchanged in structure:
HIJING~\cite{hijing} (initial conditions, strings) $\to$ string melting $\to$
ZPC parton cascade~\cite{zpc} $\to$ quark coalescence $\to$ ART hadron cascade~\cite{art}.
The framework adds hooks at each stage:

\begin{center}\small
\begin{tabular}{@{}p{3.4cm}p{6.2cm}p{5.6cm}@{}}
\toprule stage & what is added & routine (file) \\ \midrule
initialisation & read all switches & \code{HQINIT} (\code{charm\_mod.f}), called from \code{main.f} \\
HIJING hard process & triggered $c\bar c$ pair, modern PDFs, nPDF weight, HIJING trigger bug fix & \code{hijing1.383\_ampt.f}, \code{hipyset1.35.f}; \code{HQXPDF}, \code{HQNMOD}, \code{HQSHWT}, \code{HQTRGG} \\
strings $\to$ partons & charm momenta taken before Lund fragmentation & \code{HQSAVE}, \code{HQPUT}, \code{HQHTOP} \\
ZPC & charm cross section, isotropic charm scattering, collision counters, local freeze-out fixes & \code{zpc.f}, \code{HQFAC}, \code{HQCOLL}, \code{HQFRZG} \\
coalescence & He--Lin coalescence with charm $r_{BM}$; charm coalescence criterion and Peterson fragmentation; $D^*/D=1$ & \code{coal\_new.f}; \code{HQFRAG}, \code{HQDSTR} (called from \code{linana.f}) \\
ART & optional charm-hadron rescattering (not used in the reference configuration) & \code{HQHADR} (called from \code{art1f.f}) \\
output & \code{ana/charm.dat} with charm records at every stage & \code{HQHTOP}, \code{HQOUTF}, \code{HQOUTH} \\
\bottomrule
\end{tabular}
\end{center}

\section{The changes}
Each subsection gives: \textbf{why} (motivation and reference), \textbf{what} (the algorithm), \textbf{where} (routine and
lines; Tables~\ref{tab:hunks}--\ref{tab:switches} give all line numbers), \textbf{switch} and its value in the reference
configuration, and \textbf{check}.

\subsection{Heavy-quark trigger and a HIJING bug fix}
\textbf{Why.} Charm is rare: about one $c\bar c$ pair per hundred minimum-bias p+p events at RHIC. Simulating charm
observables with $10^6$ untriggered Au+Au events would give too few D mesons. HIJING has a built-in trigger
(\code{IHPR2(3)=3}) that forces one hard $c\bar c$ pair into each event~\cite{hijing}. The improved heavy-flavour AMPT
of Zheng~\emph{et al.}~\cite{zheng2020} uses it without a transverse-momentum cut on the pair.\\
\textbf{What.} \sw{AMPT\_HQTRIG=1} sets \code{IHPR2(3)=3} and \code{HIPR1(10)=}\sw{AMPT\_HQPTMIN} (default 0). One
primary nucleon--nucleon collision, chosen at random, gets a $q\bar q, gg\to c\bar c$ hard scattering (PYTHIA
\code{MSUB} 81/82). The number of primary NN collisions of the event (HIJING \code{NCOLT}) is written to the event
record, because the analysis weights each triggered event by it (Sec.~\ref{sec:analysis}).\\
\textbf{Bug fix.} \code{JETINI} (\code{hijing1.383\_ampt.f}) initialises PYTHIA lazily per collision type. Its
first-initialisation branch returned without updating the remembered type \code{ilast}. The next untriggered minijet
call then skipped the reset of \code{MSUB}/\code{CKIN}, so every later minijet of the event was generated with the
trigger process list: test events carried 24--92 charm quarks instead of about 2. The fix sets \code{ilast=itrig}
before that \code{RETURN} (tag \code{cdmeson-fix}). It is active only when the trigger is on; public AMPT does not use
the trigger.\\
\textbf{Where.} \code{HQINIT}; \code{hijing1.383\_ampt.f} hunks at the NN-collision count and \code{JETINI}
(Table~\ref{tab:hunks}). \textbf{Reference:} \sw{AMPT\_HQTRIG=1}, \sw{AMPT\_HQPTMIN} unset (0).\\
\textbf{Check.} Net charm is zero in every event (analysis counter \code{netc\_nonzero\_events}=0 on 2.4M events).

\subsection{One clean pair per event}
\textbf{Why.} With the trigger, untriggered minijets and parton showers can still produce extra charm. That breaks the
one-pair-per-event bookkeeping on which the event weight rests.\\
\textbf{What.} \sw{AMPT\_HQCLEAN=1} removes \code{MSUB} 81/82 from the untriggered minijets and sets
\code{MSTJ(45)=3} (no $g\to c\bar c$ in final-state showers).
\textbf{Where.} \code{HQINIT}; \code{hijing1.383\_ampt.f} (\code{AMPT\_HQCLEAN} hunk). \textbf{Reference:} on.

\subsection{Modern proton PDFs for the triggered pair}
\textbf{Why.} HIJING~1.383 uses the Duke--Owens PDFs at leading order. With them the $c\bar c$ cross section at
200~GeV is 6.6 times below STAR's, and the D-meson spectrum shape is distorted. Zhang~\emph{et al.}~\cite{zhang2019}
and Zheng~\emph{et al.}~\cite{zheng2020} replaced them with CTEQ6.1M in AMPT.\\
\textbf{What.} With \sw{AMPT\_HQPDF}$\geq1$, \code{PYSTFU} returns CTEQ6.1M~\cite{cteq61} (via LHAPDF~6~\cite{lhapdf},
set \code{cteq61}, member 0) for the \emph{triggered hard process only}; minijets and the bulk keep HIJING's PDFs.
The factorisation scale is $\mu_F^2=$\sw{AMPT\_HQMUF}$\times Q^2$ (default 2, i.e.\ $2m_T^2$ as in
Ref.~\cite{zheng2020}). PYTHIA's maximum search is raised by \sw{AMPT\_HQXSF} (default 1.5) so that the
acceptance--rejection step stays valid with the new, larger PDFs.\\
\textbf{Where.} \code{HQXPDF} (\code{charm\_mod.f}); hooks in \code{hipyset1.35.f} (\code{PYSTFU}, \code{PYSIGH},
\code{PYMAXI}) and \code{hijing1.383\_ampt.f} (\code{JETINI}). \textbf{Reference:} \sw{AMPT\_HQPDF=1}.\\
\textbf{Check.} The p+p D$^0$ spectrum shape against STAR gives $\chi^2/\text{ndf}=3.5/5$, against 6.7/5 for public
AMPT (\code{observation/results/tables/PP\_SHAPE.txt}).

\subsection{Impact-parameter-dependent nuclear PDFs (EPS09s) as an event weight}
\textbf{Why.} This is the change behind the main result of the paper. The charm pair is produced at a point in the
transverse plane, and its nuclear modification depends on the transverse position of the colliding nucleons in
their nuclei. Position-dependent shadowing therefore changes \emph{where} charm is produced relative to the bulk,
and hence the charm rapidity-odd ``tilt'' that drives D$^0$ directed flow. EPS09s~\cite{eps09s} is the global nPDF
fit with a spatial dependence; we use the authors' code and tables unmodified.\\
\textbf{What.} The pair is sampled with free-nucleon CTEQ6.1M PDFs. Its nuclear modification
$w=R_1(x_1,Q,s_1)\,R_2(x_2,Q,s_2)$ is then carried as an event weight (\sw{AMPT\_HQSHW=1}). $R$ is the EPS09s NLO
bound-to-free ratio of the incoming parton flavour, with isospin conjugation for neutrons, at the transverse distance
$s$ of each nucleon from its own nucleus centre (Ref.~\cite{eps09s}, Eq.~2.9). The weight is written to field 18
of the \code{charm.dat} event record and applied in the analysis. Sampling plus weighting is equivalent to sampling
with nuclear PDFs (gate g4 of the development log), keeps the events identical between nPDF variants, and lets the
31 EPS09s error sets be applied offline to the same events.\\
\textbf{Where.} \code{HQNMOD}, \code{HQSHWT}, \code{HQTRGG} (\code{charm\_mod.f}); \code{eps09s\_f.cpp} (C bindings to
the official \code{eps09s.cpp}); hooks in \code{PYSTFU}/\code{PYRAND} (\code{hipyset1.35.f}).
\textbf{Reference:} \sw{AMPT\_HQSHW=1} with \sw{AMPT\_HQPDF=1}. A related option, \sw{AMPT\_HQSHAD=s}, scales HIJING's
own shadowing of the triggered process ($R\to1+s(R-1)$). It was used only for the shadowing scan of the paper.

\subsection{Charm momentum taken before string fragmentation (direct charm)}
\textbf{Why.} In string melting, charm quarks are produced from hadrons made by Lund fragmentation, which softens
them. Zhang~\emph{et al.}~\cite{zhang2023,zhang2024} and Fan~\emph{et al.}~\cite{fan2026} take the charm quarks directly
from the HIJING hard scattering. Their spectrum is much harder.\\
\textbf{What.} \sw{AMPT\_HQDIRECT=1}. Before \code{HIJFRG}, \code{HQSAVE} stores every $c/\bar c$ on each string. After
string melting, \code{HQHTOP} gives each melted charm quark the momentum of the closest saved charm parton of the same
flavour on the same string, with formation time $0.197\,E/m_T^2$~fm. The 3-momentum difference goes to the light
partons of that string, energy-weighted, so string 3-momentum is conserved exactly. The event's energy change (a few
GeV out of $\sim40$~TeV) and the matching counters are written to the event record.\\
\textbf{Where.} \code{HQSAVE}, \code{HQPUT}, \code{HQHTOP}; hooks in \code{hijing1.383\_ampt.f} and \code{linana.f}.
\textbf{Reference:} on. \textbf{Check.} Matching failures \code{nfail}=0 in every event of the production.

\subsection{Charm--parton cross section and isotropic charm scattering in ZPC}
\textbf{Why.} With the stock forward-peaked cross section, charm couples too weakly to the medium: D$^0$ $v_2\approx0.04$
against 0.15 in data. The improved AMPT uses separate light- and heavy-quark cross sections~\cite{zhang2024}.\\
\textbf{What.} \sw{AMPT\_SIGHQ=r} multiplies the closest-approach criterion $d^2<\sigma/\pi$ by $r$ for every pair
containing a $c$ or $\bar c$, in all twelve \code{isco*} routines (\code{HQFAC}). \sw{AMPT\_HQISO=1} makes charm
collisions isotropic in the pair c.m.\ frame (same random deviate, so the event sequence is preserved), for pair c.m.\
momenta below \sw{AMPT\_HQISO\_PCUT}.\\
\textbf{Reference.} The input screening mass is 1.6016~fm$^{-1}$, i.e.\ $\sigma_{LQ}=6$~mb, and \sw{AMPT\_SIGHQ=1.5},
i.e.\ $\sigma_{HQ}=9$~mb, isotropic below 0.5~GeV. These values were chosen jointly on D$^0$ $v_2$ and $R_{AA}$.

\subsection{Local parton freeze-out}
\textbf{Why.} The reference input uses AMPT's \code{ISOFT=5} option (``local freeze-out'' in Lin's \code{linana.f}): a
parton freezes when its local energy density falls below a critical value, rather than after its last collision. In
the stock code, a parton that never satisfies the criterion before the cascade ends is hadronised with its
\emph{initial} momentum. We fix that.\\
\textbf{What.} \sw{AMPT\_ECRIT} sets the critical energy density (stock 1.0~GeV/fm$^3$). \code{HQFRZG} freezes a
never-frozen parton at its last interaction, as \code{ISOFT=4} does (0.07\,\% of partons are affected). With
\sw{AMPT\_HQLATE=0}, charm quarks never freeze by the density criterion. They keep scattering, changing only their own
momentum, until the cascade ends, and are then frozen at their last interaction.\\
\textbf{Where.} \code{amptsub.f} (\code{ARTSET}), \code{HQFRZG}, \code{zpc.f}, \code{linana.f} \code{LOCAL}.
\textbf{Reference:} \code{ISOFT=5}, \sw{AMPT\_ECRIT=1.0}, \sw{AMPT\_HQLATE=0}.

\subsection{Hadronisation: improved coalescence, charm coalescence criterion and fragmentation}
\textbf{Why.} The public AMPT coalescence forms mesons and baryons from fixed quark groupings and builds antibaryons from
antiquarks a median 13~fm apart. He and Lin's improved coalescence~\cite{helin2017} chooses the closest partners and
lets a parameter $r_{BM}$ decide between baryon and meson formation; it is published but not distributed, so we
reimplemented it. For charm, Zheng~\emph{et al.}~\cite{zheng2020} use a larger $r_{BM}$ to produce enough $\Lambda_c$.
Zhang~\emph{et al.}~\cite{zhang2024} and Fan~\emph{et al.}~\cite{fan2026} add a coalescence criterion: charm hadrons
whose constituents are too far apart, too massive or at too high $m_T$ fragment instead (Peterson
fragmentation~\cite{peterson}).\\
\textbf{What.}
\begin{itemize}
\item \sw{AMPT\_COALNEW=1}: He--Lin coalescence (\code{coal\_new.f}) with $r_{BM}=0.61$ for light flavours and
  $r_{BM}^c=$\sw{AMPT\_RBMHQ} for groups containing charm. \sw{AMPT\_CNORDER=2} sweeps the partons in a random order
  per event, removing a projectile/target bias of the stored order.
\item \sw{AMPT\_HQFRAG=1} (\code{HQFRAG}): an open-charm candidate coalesces if $m_T<m_T^{\max}$ (\sw{AMPT\_HQMTMAX}),
  or if its largest constituent separation is below $p_r$ (\sw{AMPT\_HQPR}) and $m_{\rm inv}<\sum m_q+p_m(m_H-\sum m_q)$
  (\sw{AMPT\_HQPM}), with constituent quark masses (\sw{AMPT\_HQCONST=1}, Fan~\emph{et al.} Eq.~2) (\sw{AMPT\_HQMTMODE=1}:
  the ``or'' of Ref.~\cite{fan2026}). Otherwise the charm quark fragments. $z$ is drawn from the Peterson function
  ($\epsilon_c=0.05$), the light antiquark comes from a vacuum pair ($u:d:s=1:1:0.3$), and the remaining momentum goes
  to a light hadron, so 3-momentum is conserved per candidate. The hadron record carries a mode flag
  (1 coalescence, 2 fragmentation). \sw{AMPT\_HQBDIST=1} (systematic variant c1) uses the Jacobi distance
  $\sqrt{2/3}\,|r_c-(r_{q_1}+r_{q_2})/2|$ for baryons.
\item \sw{AMPT\_DSTAR1=1} (\code{HQDSTR}): primordial $D^*/D=1$ (and $D_s^*/D_s=1$), as in Ref.~\cite{zheng2020}. The
  charm--light candidates are ordered by excess mass; the lower half become $D$, the upper half $D^*$.
\end{itemize}
\textbf{Reference:} \sw{AMPT\_COALNEW=1}, \sw{AMPT\_RBMHQ=3}, \sw{AMPT\_CNORDER=2}, \sw{AMPT\_HQFRAG=1}, $p_r=0.5$~fm,
$p_m=0.5$, $m_T^{\max}=3.3$ (Au+Au) and 2.0~GeV (p+p), \sw{AMPT\_HQCONST=1}, \sw{AMPT\_DSTAR1=1}.

\subsection{Output}
\sw{AMPT\_CHARM\_OUT=1} writes \code{ana/charm.dat} with one record type per stage. E: event, with $b$, the
participants, $N_{\rm coll}$, the energy change, the triggered $\sigma_{c\bar c}$ and the nPDF weight. I/L: charm and
a light-parton sample at formation. F/G: the same at freeze-out. H: charm hadrons after coalescence, with mode and
constituent separation. Z: the end of the parton phase. A private random generator is used, so writing the file never
changes the events. \sw{AMPT\_LQFRAC} sets the light-parton sampling fraction. The node-side reduction
\code{scripts/reduce\_dmeson\_v2.awk} keeps every charm hadron, all charged kaons and protons, a 25\,\% (10\,\%
forward) sample of pions with weights, and the full second-order flow vectors of all charged hadrons in the STAR and
PHENIX event-plane windows.

\subsection{Options implemented but not used in the reference configuration}
\begin{itemize}
\item \sw{AMPT\_HQHAD=1} (\code{HQHADR}): elastic, isotropic charm--hadron rescattering in ART with the He--Fries--Rapp
  cross sections (7~mb with mesons, 10~mb with baryons)~\cite{hfr2011}. It was tested on the reference configuration:
  0.17 collisions per charm hadron, $v_2$ unchanged, $R_{AA}$ worse by 4--14\,\%. Not adopted; quoted as a
  systematic in the paper.
\item \sw{AMPT\_HQCRONIN=$\delta$} (\code{HQCKICK}): Cronin $k_T$ broadening of the pair, Zhang~\emph{et
  al.}~\cite{zhang2023} Eqs.~1--3 (used by Fan~\emph{et al.}; a candidate improvement for $R_{AA}$).
\item \sw{AMPT\_ZPC\_TSTOP}: hard stop of the parton cascade, Zhang and Xu~\cite{zhangxu2017}.
\item \sw{AMPT\_TILT}: a sensitivity knob (light partons shifted by $\lambda\eta_s$), not a physics model.
\item \sw{AMPT\_SOUT}, \sw{AMPT\_HQDREST}, \sw{AMPT\_HQFBAR}, \sw{AMPT\_HQPDFSET}, \sw{AMPT\_HQCDBG}: diagnostics and
  variants.
\end{itemize}

\section{Reference configuration}
Input card \code{inputs/auau200\_lf\_s6.ampt} = the stock \code{input.ampt} with \code{ISOFT=5} and screening mass
1.6016~fm$^{-1}$ (the event count, impact parameter range and seeds are set per job). Switches:
\code{configs/reference\_auau.env} and \code{configs/reference\_pp.env} (Table~\ref{tab:switches}). Public AMPT for
comparison: \code{inputs/auau200\_sm.ampt} (the stock card) with only the trigger and the output switches
(\code{configs/public\_*.env}).

\begin{center}\small
\begin{tabular}{@{}lllll@{}}
\toprule arm & config & events & HIJING seed base & ZPC seed \\ \midrule
dm24\_e\_a, dm24\_e\_b & reference\_auau & $2\times1000\times300$ & 20261107, 20261108 & $13i+9$, $13i+10$ \\
dm25\_e\_c\ldots h & reference\_auau & $6\times1000\times300$ & 20261109\ldots14 & $31i+11\ldots16$ \\
dm24\_epp, dm25\_epp & reference\_pp & $2\times200\times30000$ & 20261107, 20261115 & $13i+9$, $31i+17$ \\
dm25\_c1 / dm25\_ppc1 & c1\_auau / c1\_pp & $1000\times300$ / $100\times30000$ & 20261116 / 20261117 & $31i+18$ / $31i+19$ \\
dm25\_pub\_a, b / dm25\_pppub & public\_* & $2\times1000\times300$ / $100\times30000$ & 20261118, 19 / 20261120 & $31i+20, 21$ / $31i+22$ \\
\bottomrule
\end{tabular}
\end{center}
HIJING seed of job $i$ = base $+\,7919\,i$. Au+Au: $0<b<14.5$~fm; p+p: $0<b<2$~fm.

\section{Analysis conventions}\label{sec:analysis}
\begin{itemize}
\item \textbf{Feed-down.} $D^{*0}\to D^0$, $D^{*+}\to D^0\pi^+$ (67.7\,\%), $D_s^*\to D_s$, $\Sigma_c\to\Lambda_c$
  (PDG branching ratios, \code{analysis/dmlib.py}), as in the measured yields.
\item \textbf{Weights.} Each triggered Au+Au event carries one pair, while real charm production scales with
  $N_{\rm coll}$. Charm observables are therefore weighted by the event's \code{NCOLT} times the EPS09s weight; the
  charm-quark count by \code{NCOLT} only. Light hadrons are never weighted.
\item \textbf{$R_{AA}$} $=$ (D$^0$ per $c$ quark, Au+Au) / (D$^0$ per $c$ quark, p+p). This equals the usual definition
  when the total charm yield scales with $N_{\rm coll}$ times the nPDF factor.
\item \textbf{Centrality.} Charged-particle multiplicity in $|\eta|<0.5$, as STAR.
\item \textbf{Flow.} $v_2$: the scalar product in STAR's $\eta$-sub-event geometry, with the resolution taken per 5\,\%
  bin; $v_1$: relative to the reaction plane (STAR uses the spectator plane), with the sign fixed by HIJING's projectile
  at $x=+b/2$.
\item \textbf{Errors.} Jackknife over job files (delete-a-block, at most 2000 blocks).
\end{itemize}

\section{Verification}\label{sec:verify}
\begin{itemize}
\item \textbf{Stock reproduction (G0).} With no switch set, \code{ana/ampt.dat} is byte-identical to the unmodified
  program.
\item \textbf{Exact reproduction of the production} (\code{make test-repro}). The first events of Au+Au job~1 and p+p
  job~1 of the production, regenerated on a different machine and compiler (gfortran 13.3 vs 11.5), give a
  byte-identical \code{charm.dat} and identical event headers, flow vectors and charm hadrons.
\item \textbf{Independent re-analysis} (\code{make test-quick}). \code{analysis/indep\_check.py} re-implements the
  observables from the documentation without importing the production analysis. On 36k Au+Au and 300k p+p events it
  agrees with \code{analysis/analyze\_arm.py} to $10^{-17}$.
\item \textbf{Results} (\code{make test-results}). Every table and figure in \code{observation/results/} is
  regenerated from the shipped analysis outputs.
\end{itemize}

\begin{thebibliography}{99}\small
\bibitem{ampt2005} Z.-W. Lin, C. M. Ko, B.-A. Li, B. Zhang, S. Pal, Phys. Rev. C \textbf{72}, 064901 (2005).
\bibitem{lin2021} Z.-W. Lin, L. Zheng, Nucl. Sci. Tech. \textbf{32}, 113 (2021), arXiv:2110.02989.
\bibitem{hijing} X.-N. Wang, M. Gyulassy, Phys. Rev. D \textbf{44}, 3501 (1991); M. Gyulassy, X.-N. Wang, Comput.
  Phys. Commun. \textbf{83}, 307 (1994).
\bibitem{zpc} B. Zhang, Comput. Phys. Commun. \textbf{109}, 193 (1998).
\bibitem{art} B.-A. Li, C. M. Ko, Phys. Rev. C \textbf{52}, 2037 (1995).
\bibitem{zheng2020} L. Zheng, C. Zhang, S. S. Shi, Z.-W. Lin, Phys. Rev. C \textbf{101}, 034905 (2020).
\bibitem{zhang2019} C. Zhang, L. Zheng, F. Liu, S. Shi, Z.-W. Lin, Phys. Rev. C \textbf{99}, 064906 (2019).
\bibitem{zhang2023} C. Zhang, L. Zheng, S. Shi, Z.-W. Lin, Phys. Lett. B \textbf{846}, 138219 (2023).
\bibitem{zhang2024} C. Zhang, L. Zheng, S. Shi, Z.-W. Lin, Eur. Phys. J. C \textbf{84}, 942 (2024).
\bibitem{fan2026} B. Fan, S. Shi, C. Zhang, L. Zheng, arXiv:2601.12287 (2026).
\bibitem{helin2017} Y. He, Z.-W. Lin, Phys. Rev. C \textbf{96}, 014910 (2017).
\bibitem{peterson} C. Peterson, D. Schlatter, I. Schmitt, P. M. Zerwas, Phys. Rev. D \textbf{27}, 105 (1983).
\bibitem{cteq61} D. Stump \emph{et al.}, JHEP \textbf{10} (2003) 046.
\bibitem{lhapdf} A. Buckley \emph{et al.}, Eur. Phys. J. C \textbf{75}, 132 (2015).
\bibitem{eps09s} I. Helenius, K. J. Eskola, H. Honkanen, C. A. Salgado, JHEP \textbf{07} (2012) 073.
\bibitem{hfr2011} M. He, R. J. Fries, R. Rapp, Phys. Lett. B \textbf{701}, 445 (2011).
\bibitem{zhangxu2017} C.-J. Zhang, J. Xu, Phys. Rev. C \textbf{96}, 044907 (2017).
\end{thebibliography}

\appendix
\section{Code map (generated from the repository)}
{\small\begin{longtable}{@{}p{2.7cm}p{1.5cm}p{1.3cm}p{9.4cm}@{}}
\caption{Every hunk of \texttt{patches/ampt-v1.26t9b-v2.26t9b\_charm.patch}. Lines = the added lines in the patched file \texttt{build/ampt-src/}; $+/-$ = lines added/removed. The note is the first comment of the hunk (tagged \texttt{cdmeson} in the source).}\label{tab:hunks}\\
\toprule file & lines & $+/-$ & note \\ \midrule \endfirsthead
\toprule file & lines & $+/-$ & note \\ \midrule \endhead
\texttt{amptsub.f} & 44 & +1/$-$0 & code: include 'hqctl.inc' \\
\texttt{amptsub.f} & 52--53 & +2/$-$0 & cdmeson (src\_\allowbreak{}v9): AMPT\_\allowbreak{}ECRIT overrides it (read by HQINIT) \\
\texttt{amptsub.f} & 4923--4925 & +3/$-$1 & cdmeson (src\_\allowbreak{}v9): partons never frozen keep their final, not initial, \\
\texttt{art1f.f} & 619--621 & +3/$-$1 & dmeson src\_\allowbreak{}v17: charm-hadron rescattering (returns at once unless \\
\texttt{hijing1.383\_\allowbreak{}ampt.f} & 199--200 & +2/$-$0 & cdmeson src\_\allowbreak{}v13: ihqtrx (trigger context of AMPT\_\allowbreak{}HQSHAD) \\
\texttt{hijing1.383\_\allowbreak{}ampt.f} & 583--599 & +12/$-$0 & cdmeson: number of primary NN collisions of this event (E record) \\
\texttt{hijing1.383\_\allowbreak{}ampt.f} & 648--655 & +6/$-$0 & cdmeson src\_\allowbreak{}v14: with AMPT\_\allowbreak{}HQPDF the triggered PYINIT (maximum search) \\
\texttt{hijing1.383\_\allowbreak{}ampt.f} & 1346--1349 & +3/$-$0 & cdmeson: record charm partons on the strings before Lund fragmentation \\
\texttt{hijing1.383\_\allowbreak{}ampt.f} & 4298--4299 & +2/$-$0 & cdmeson: \\
\texttt{hijing1.383\_\allowbreak{}ampt.f} & 4396--4400 & +5/$-$0 & cdmeson: AMPT\_\allowbreak{}HQCLEAN, no heavy-quark pair production in untriggered jets \\
\texttt{hijing1.383\_\allowbreak{}ampt.f} & 4527--4547 & +18/$-$1 & cdmeson src\_\allowbreak{}v14: AMPT\_\allowbreak{}HQPDF raises every triggered maximum by hqxsf: \\
\texttt{hipyset1.35.f} & 7603--7604 & +2/$-$0 & cdmeson src\_\allowbreak{}v14: \\
\texttt{hipyset1.35.f} & 7874--7875 & +2/$-$0 & cdmeson src\_\allowbreak{}v14: largest weight/maximum inside the triggered process \\
\texttt{hipyset1.35.f} & 7970--7971 & +2/$-$0 & cdmeson src\_\allowbreak{}v15: AMPT\_\allowbreak{}HQSHW=1, nuclear modification as event weight \\
\texttt{hipyset1.35.f} & 12017--12018 & +2/$-$0 & cdmeson src\_\allowbreak{}v14: \\
\texttt{hipyset1.35.f} & 12144--12147 & +3/$-$0 & cdmeson src\_\allowbreak{}v14: hard-process call (AMPT\_\allowbreak{}HQMUF scales muF\^\{\}2 here only) \\
\texttt{hipyset1.35.f} & 13981--13982 & +2/$-$0 & cdmeson src\_\allowbreak{}v13: \\
\texttt{hipyset1.35.f} & 14331--14336 & +6/$-$0 & cdmeson src\_\allowbreak{}v14: AMPT\_\allowbreak{}HQPDF, triggered hard process: CTEQ6.1M (x EPS09s \\
\texttt{hipyset1.35.f} & 14575--14578 & +4/$-$0 & cdmeson src\_\allowbreak{}v14: AMPT\_\allowbreak{}HQPDF already applied EPS09s (no HIJING shadowing) \\
\texttt{hipyset1.35.f} & 14597--14598 & +2/$-$1 & cdmeson src\_\allowbreak{}v13: AMPT\_\allowbreak{}HQSHAD scales the triggered hard process only \\
\texttt{linana.f} & 656--657 & +2/$-$0 & cdmeson: direct-charm option, ZPC charm counters, charm.dat records \\
\texttt{linana.f} & 734--738 & +5/$-$1 & cdmeson: charm partons at ZPC freeze-out \\
\texttt{linana.f} & 929--930 & +2/$-$0 & cdmeson: AMPT\_\allowbreak{}DSTAR1, primordial D*/D = 1 (Zheng et al. 2020) \\
\texttt{linana.f} & 1099--1102 & +4/$-$0 & cdmeson: AMPT\_\allowbreak{}HQFRAG, coalescence criterion + Peterson fragmentation \\
\texttt{linana.f} & 1998 & +1/$-$0 & code: include 'hqctl.inc' \\
\texttt{linana.f} & 2047--2057 & +11/$-$2 & cdmeson src\_\allowbreak{}v12b (AMPT\_\allowbreak{}HQLATE=x>=0): a charm quark freezes only below \\
\texttt{main.f} & 179--180 & +2/$-$1 & cdmeson: D-meson framework switches (charm\_\allowbreak{}mod.f), all default off \\
\texttt{zpc.f} & 1121 & +1/$-$0 & code: include 'hqctl.inc' \\
\texttt{zpc.f} & 1134--1140 & +7/$-$0 & cdmeson (src\_\allowbreak{}v9): AMPT\_\allowbreak{}ZPC\_\allowbreak{}TSTOP, hard stop of the partonic cascade at \\
\texttt{zpc.f} & 1864--1865 & +2/$-$0 & cdmeson: count charm scatterings \\
\texttt{zpc.f} & 1972--2001 & +24/$-$3 & cdmeson src\_\allowbreak{}v12 (AMPT\_\allowbreak{}HQLATE>=0): an unfrozen charm quark still scatters \\
\texttt{zpc.f} & 2196--2207 & +12/$-$0 & cdmeson src\_\allowbreak{}v12: the frozen partner of a late charm collision is untouched \\
\texttt{zpc.f} & 2223 & +1/$-$0 & code: logical hqisof \\
\texttt{zpc.f} & 2239--2240 & +2/$-$0 & cdmeson: AMPT\_\allowbreak{}HQISO, isotropic scattering for charm (same deviate rx) \\
\texttt{zpc.f} & 3936 & +1/$-$1 & code: allok = allok .and. dm2 .lt. cutof2*hqfac(i, j) \\
\texttt{zpc.f} & 4064 & +1/$-$1 & code: allok = allok .and. dm2 .lt. cutof2*hqfac(i, j) \\
\texttt{zpc.f} & 4199 & +1/$-$1 & code: allok = allok .and. dm2 .lt. cutof2*hqfac(i, j) \\
\texttt{zpc.f} & 4349 & +1/$-$1 & code: allok = allok .and. dm2 .lt. cutof2*hqfac(i, j) \\
\texttt{zpc.f} & 4505 & +1/$-$1 & code: allok = allok .and. dm2 .lt. cutof2*hqfac(i, j) \\
\texttt{zpc.f} & 4669 & +1/$-$1 & code: allok = allok .and. dm2 .lt. cutof2*hqfac(i, j) \\
\texttt{zpc.f} & 4802 & +1/$-$1 & code: allokp = allokp .and. dm2 .lt. cutof2*hqfac(i, j) \\
\texttt{zpc.f} & 4966 & +1/$-$1 & code: allokp = allokp .and. dm2 .lt. cutof2*hqfac(i, j) \\
\texttt{zpc.f} & 5146 & +1/$-$1 & code: allokp = allokp .and. dm2 .lt. cutof2*hqfac(i, j) \\
\texttt{zpc.f} & 5294 & +1/$-$1 & code: allokp = allokp .and. dm2 .lt. cutof2*hqfac(i, j) \\
\texttt{zpc.f} & 5462 & +1/$-$1 & code: allokp = allokp .and. dm2 .lt. cutof2*hqfac(i, j) \\
\texttt{zpc.f} & 5646 & +1/$-$1 & code: allokp = allokp .and. dm2 .lt. cutof2*hqfac(i, j) \\
\bottomrule \end{longtable}}

{\small\begin{longtable}{@{}p{1.9cm}p{2.0cm}p{1.6cm}p{9.4cm}@{}}
\caption{Routines of the new files (\texttt{src\_new/}), with their line ranges and the first line of their header comment.}\label{tab:routines}\\
\toprule file & routine & lines & purpose \\ \midrule \endfirsthead
\toprule file & routine & lines & purpose \\ \midrule \endhead
\texttt{charm\_mod.f} & \texttt{HQBLK} & 72--92 &  \\
\texttt{charm\_mod.f} & \texttt{HQINIT} & 93--258 & read the switches; called from main.f before HIJSET \\
\texttt{charm\_mod.f} & \texttt{HQFAC} & 259--273 & ZPC: factor on cutof2 (=sigma/pi) for the pair (i,j) in /prec2/ \\
\texttt{charm\_mod.f} & \texttt{HQCOLL} & 274--306 & ZPC: count 2-body scatterings of charm partons (/prec2/ index) \\
\texttt{charm\_mod.f} & \texttt{HQFRZG} & 307--348 & ISOFT=5, end of ZPC (zpstrg, before /frzprc/ replaces /prec2/): \\
\texttt{charm\_mod.f} & \texttt{HQSAVE} & 349--398 & HIJING, string-melting branch.  ISG=0: reset for this event and \\
\texttt{charm\_mod.f} & \texttt{HQPUT} & 399--419 &  \\
\texttt{charm\_mod.f} & \texttt{HQHTOP} & 420--592 & end of HTOP: (a) AMPT\_\allowbreak{}HQDIRECT replacement, (b) reset ZPC charm \\
\texttt{charm\_mod.f} & \texttt{hqrand} & 593--605 & private LCG for output subsampling; never touches RANART/NSEED so \\
\texttt{charm\_mod.f} & \texttt{HQOUTF} & 606--646 & start of PTOH (ZPC finished): charm partons at freeze-out (F) with \\
\texttt{charm\_mod.f} & \texttt{HQOUTH} & 647--715 & end of PTOH: charm hadrons right after coalescence (H) with the \\
\texttt{charm\_mod.f} & \texttt{HQNCOL} & 716--723 & HIJING: store the number of primary NN collisions of the event \\
\texttt{charm\_mod.f} & \texttt{HQDSTR} & 724--827 & AMPT\_\allowbreak{}DSTAR1 (PTOH, right after the flavour/species loop): set the \\
\texttt{charm\_mod.f} & \texttt{HQFRAG} & 828--1156 & end of PTOH (hadron momenta/positions set): coalescence criterion \\
\texttt{charm\_mod.f} & \texttt{KBARY} & 1157--1186 & lightest ground-state baryon (KF > 0) of three light flavours \\
\texttt{charm\_mod.f} & \texttt{KCBARY} & 1187--1205 & ground-state singly-charmed baryon (KF > 0) c + light diquark \\
\texttt{charm\_mod.f} & \texttt{HQXPDF} & 1206--1248 & PYSTFU, triggered hard process only: x*f of a proton from CTEQ6.1M \\
\texttt{charm\_mod.f} & \texttt{HQNMOD} & 1249--1299 & src\_\allowbreak{}v15 (moved out of HQXPDF, same arithmetic): EPS09s NLO bound- \\
\texttt{charm\_mod.f} & \texttt{HQSHWT} & 1300--1353 & src\_\allowbreak{}v15 AMPT\_\allowbreak{}HQSHW=1, PYRAND after the channel choice of the \\
\texttt{charm\_mod.f} & \texttt{HQRRX} & 1354--1383 & src\_\allowbreak{}v15: HIJING 1.383 nuclear shadowing factor exactly as PYSTFU \\
\texttt{charm\_mod.f} & \texttt{hqrnp} & 1384--1391 & EPS09s power series 1 + c1 T + c2 T\^\{\}2 + c3 T\^\{\}3 + c4 T\^\{\}4 \\
\texttt{charm\_mod.f} & \texttt{HQTRGG} & 1392--1409 & HIJING, at the choice of the triggered NN collision: number of \\
\texttt{charm\_mod.f} & \texttt{HQCKICK} & 1410--1569 & AMPT\_\allowbreak{}HQCRONIN=delta, end of HTOP after the HQDIRECT replacement: \\
\texttt{charm\_mod.f} & \texttt{HQLBST} & 1570--1585 & Lorentz boost of the four-vector q (px,py,pz,E) by velocity b \\
\texttt{charm\_mod.f} & \texttt{HQVMAX} & 1586--1602 & once per event: report a new maximum of SIGS/XSEC\_\allowbreak{}max inside the \\
\texttt{charm\_mod.f} & \texttt{HQHADR} & 1603--1815 & src\_\allowbreak{}v17 AMPT\_\allowbreak{}HQHAD=1: hadronic rescattering of charm hadrons in \\
\texttt{charm\_mod.f} & \texttt{HQBST4} & 1816--1835 & src\_\allowbreak{}v17: Lorentz boost of the 4-vector x = (x1,x2,x3,x0) into the \\
\texttt{charm\_mod.f} & \texttt{hqrnh} & 1836--1846 & src\_\allowbreak{}v17: private Park-Miller generator of the charm rescattering \\
\texttt{coal\_new.f} & \texttt{CNPAR} & 1--72 & He \& Lin's improved quark coalescence, PRC 96 014910 (arXiv:1703.02673), \\
\texttt{coal\_new.f} & \texttt{CNDIS2} & 73--115 & Relative distance between two flat-list partons in AMPT's own convention: \\
\texttt{coal\_new.f} & \texttt{CNDIS3} & 116--163 & He \& Lin's d\_\allowbreak{}B: the average of the three relative distances among three \\
\texttt{coal\_new.f} & \texttt{COALNEW} & 164--366 & Replaces `coales` when AMPT\_\allowbreak{}COALNEW=1. \\
\texttt{coal\_new.f} & \texttt{CNCHK} & 367--462 & Gate 1: parton bookkeeping, as an assertion. A hadronisation rewrite that \\
\texttt{coal\_new.f} & \texttt{CNWRIT} & 463--527 & Write the chosen groups back into the SOFT arrays and reset NSG, so that \\
\bottomrule \end{longtable}}

{\small\begin{longtable}{@{}p{4.2cm}p{3.4cm}p{2.8cm}p{2.8cm}@{}}
\caption{Environment switches: the source line that reads each one, and its value in the reference Au+Au and p+p configurations (\texttt{configs/reference\_*.env}). Unset = the stock behaviour.}\label{tab:switches}\\
\toprule switch & read at & reference Au+Au & reference p+p \\ \midrule \endfirsthead
\toprule switch & read at & reference Au+Au & reference p+p \\ \midrule \endhead
\texttt{AMPT\_\allowbreak{}HQTRIG} & \texttt{charm\_\allowbreak{}mod.f:100} & 1 & 1 \\
\texttt{AMPT\_\allowbreak{}HQCLEAN} & \texttt{charm\_\allowbreak{}mod.f:102} & 1 & 1 \\
\texttt{AMPT\_\allowbreak{}HQDIRECT} & \texttt{charm\_\allowbreak{}mod.f:104} & 1 & 1 \\
\texttt{AMPT\_\allowbreak{}CHARM\_\allowbreak{}OUT} & \texttt{charm\_\allowbreak{}mod.f:106} & 1 & 1 \\
\texttt{AMPT\_\allowbreak{}DSTAR1} & \texttt{charm\_\allowbreak{}mod.f:108} & 1 & 1 \\
\texttt{AMPT\_\allowbreak{}HQISO} & \texttt{charm\_\allowbreak{}mod.f:110} & 1 & unset \\
\texttt{AMPT\_\allowbreak{}HQISO\_\allowbreak{}PCUT} & \texttt{charm\_\allowbreak{}mod.f:112} & 0.5 & unset \\
\texttt{AMPT\_\allowbreak{}TILT} & \texttt{charm\_\allowbreak{}mod.f:114} & unset & unset \\
\texttt{AMPT\_\allowbreak{}ECRIT} & \texttt{charm\_\allowbreak{}mod.f:119} & 1.0 & 1.0 \\
\texttt{AMPT\_\allowbreak{}SOUT} & \texttt{charm\_\allowbreak{}mod.f:124} & unset & unset \\
\texttt{AMPT\_\allowbreak{}HQLATE} & \texttt{charm\_\allowbreak{}mod.f:126} & 0 & 0 \\
\texttt{AMPT\_\allowbreak{}HQSHAD} & \texttt{charm\_\allowbreak{}mod.f:132} & unset & unset \\
\texttt{AMPT\_\allowbreak{}HQPDF} & \texttt{charm\_\allowbreak{}mod.f:138} & 1 & 1 \\
\texttt{AMPT\_\allowbreak{}HQPDFSET} & \texttt{charm\_\allowbreak{}mod.f:140} & unset & unset \\
\texttt{AMPT\_\allowbreak{}HQMUF} & \texttt{charm\_\allowbreak{}mod.f:142} & unset & unset \\
\texttt{AMPT\_\allowbreak{}HQXSF} & \texttt{charm\_\allowbreak{}mod.f:144} & unset & unset \\
\texttt{AMPT\_\allowbreak{}HQDREST} & \texttt{charm\_\allowbreak{}mod.f:158} & unset & unset \\
\texttt{AMPT\_\allowbreak{}HQSHW} & \texttt{charm\_\allowbreak{}mod.f:163} & 1 & unset \\
\texttt{AMPT\_\allowbreak{}HQCONST} & \texttt{charm\_\allowbreak{}mod.f:168} & 1 & 1 \\
\texttt{AMPT\_\allowbreak{}HQBDIST} & \texttt{charm\_\allowbreak{}mod.f:174} & unset & unset \\
\texttt{AMPT\_\allowbreak{}HQHAD} & \texttt{charm\_\allowbreak{}mod.f:177} & unset & unset \\
\texttt{AMPT\_\allowbreak{}HQHADM} & \texttt{charm\_\allowbreak{}mod.f:179} & unset & unset \\
\texttt{AMPT\_\allowbreak{}HQHADB} & \texttt{charm\_\allowbreak{}mod.f:181} & unset & unset \\
\texttt{AMPT\_\allowbreak{}HQHADDS} & \texttt{charm\_\allowbreak{}mod.f:183} & unset & unset \\
\texttt{AMPT\_\allowbreak{}HQHADLC} & \texttt{charm\_\allowbreak{}mod.f:185} & unset & unset \\
\texttt{AMPT\_\allowbreak{}HQHADREC} & \texttt{charm\_\allowbreak{}mod.f:187} & unset & unset \\
\texttt{AMPT\_\allowbreak{}HQCRONIN} & \texttt{charm\_\allowbreak{}mod.f:197} & unset & unset \\
\texttt{AMPT\_\allowbreak{}ZPC\_\allowbreak{}TSTOP} & \texttt{charm\_\allowbreak{}mod.f:202} & unset & unset \\
\texttt{AMPT\_\allowbreak{}HQMTMODE} & \texttt{charm\_\allowbreak{}mod.f:207} & 1 & 1 \\
\texttt{AMPT\_\allowbreak{}HQFBAR} & \texttt{charm\_\allowbreak{}mod.f:209} & unset & unset \\
\texttt{AMPT\_\allowbreak{}HQFRAG} & \texttt{charm\_\allowbreak{}mod.f:211} & 1 & 1 \\
\texttt{AMPT\_\allowbreak{}HQPR} & \texttt{charm\_\allowbreak{}mod.f:213} & 0.5 & 0.5 \\
\texttt{AMPT\_\allowbreak{}HQPM} & \texttt{charm\_\allowbreak{}mod.f:215} & 0.5 & 0.5 \\
\texttt{AMPT\_\allowbreak{}HQMTMAX} & \texttt{charm\_\allowbreak{}mod.f:217} & 3.3 & 2.0 \\
\texttt{AMPT\_\allowbreak{}HQEPS} & \texttt{charm\_\allowbreak{}mod.f:219} & unset & unset \\
\texttt{AMPT\_\allowbreak{}HQPTMIN} & \texttt{charm\_\allowbreak{}mod.f:221} & unset & unset \\
\texttt{AMPT\_\allowbreak{}SIGHQ} & \texttt{charm\_\allowbreak{}mod.f:223} & 1.5 & 1.5 \\
\texttt{AMPT\_\allowbreak{}LQFRAC} & \texttt{charm\_\allowbreak{}mod.f:225} & 0.05 & unset \\
\texttt{AMPT\_\allowbreak{}HQCDBG} & \texttt{charm\_\allowbreak{}mod.f:879} & unset & unset \\
\texttt{AMPT\_\allowbreak{}COALNEW} & \texttt{coal\_\allowbreak{}new.f:24} & 1 & 1 \\
\texttt{AMPT\_\allowbreak{}CNORDER} & \texttt{coal\_\allowbreak{}new.f:40} & 2 & 2 \\
\texttt{AMPT\_\allowbreak{}RBM} & \texttt{coal\_\allowbreak{}new.f:47} & unset & unset \\
\texttt{AMPT\_\allowbreak{}RBMHQ} & \texttt{coal\_\allowbreak{}new.f:57} & 3 & 3 \\
\bottomrule \end{longtable}}

\end{document}
